\section{One free dial}
\label{sec:ssrdial}

We know the sign conventions; now build the rule.  The guiding idea is
modest: for small moves, a transport should not invent new smile
\emph{shape} --- no quotes have printed that would justify a new
shape.  A curve can be changed in only two ways without changing its
shape: slide its argument, or add a constant.  A shape-preserving
transport is therefore
\[
\sigma_{\mathrm{new}}(k)
=\sigma_{\mathrm{old}}(k+aH)+bH
\]
for two numbers $a$ and $b$ per unit move.  The slide is fixed by
anchoring the family to the one hypothesis whose meaning is
operational: if the level term is zero, the transport should be
exactly sticky-strike --- volatilities pinned to strikes, the pure
re-indexing \eqref{eq:ssrcurverule} of the last section.  That sets
$a=1$.  What remains free is one number: the level.

Write $s_0$ for the \emph{ATM skew}: the slope of the smile at the
money, $s_0=\partial_k\sigma|_{k=0}$, in volatility per unit of
log-moneyness.  On the frozen hero node $s_0=\MacSsrHeroSkew$: the
fitted curve loses about four volatility points per ten percent of
moneyness as $k$ rises through zero.  The natural yardstick for the
level is then $s_0H$ --- the skew times the move, the amount of
volatility the smile's own tilt attaches to a move of size $H$.
Writing the level as $(\mathcal{R}-1)\,s_0H$, a multiple of that
yardstick --- the reason for the $-1$ appears in one moment --- gives
the chapter's transport, its one boxed display:
\begin{equation}
\boxed{\;\sigma_{\mathrm{new}}(k)
=\sigma_{\mathrm{old}}(k+H)+(\mathcal{R}-1)\,s_0\,H .\;}
\label{eq:ssrshift}
\end{equation}

The dial $\mathcal{R}$ takes its meaning at the money.

\begin{proposition}[ATM response]
\label{prop:ssratm}
Under \eqref{eq:ssrshift}, the at-the-money volatility responds to the
move by
\[
\sigma_{\mathrm{new}}(0)-\sigma_{\mathrm{old}}(0)
=\mathcal{R}\,s_0\,H+O(H^{2}) ,
\qquad\text{hence}\qquad
\mathcal{R}
=\frac{\dd\sigma_{\mathrm{atm}}/\dd\log F}{s_0} .
\]
\end{proposition}

\begin{proof}
Evaluate \eqref{eq:ssrshift} at $k=0$ and expand the re-indexed term
to first order in $H$:
\[
\sigma_{\mathrm{new}}(0)
=\sigma_{\mathrm{old}}(H)+(\mathcal{R}-1)s_0H
=\sigma_{\mathrm{old}}(0)+s_0H+(\mathcal{R}-1)s_0H+O(H^{2})
=\sigma_{\mathrm{old}}(0)+\mathcal{R}s_0H+O(H^{2}) .
\]
The second equality is the point: the re-indexing alone already moves
the ATM reading by $s_0H$, because sliding along a tilted curve
changes the value at the money.  The level term adds the rest.
\end{proof}

So $\mathcal{R}$ is the fraction of the skew realized as ATM movement
per unit log-forward move.  The literature calls this the
\emph{skew-stickiness ratio} \citep{Bergomi2016}, and the classical
taxonomy of \citet{Derman1999} names three of its values.  Read them
through a falling market, since $s_0<0$: the forward drops
($H<0$), and the question is what the ATM volatility does.
\begin{itemize}
\item $\mathcal{R}=0$, \emph{sticky-moneyness}: to first order the
ATM volatility does not move; the whole smile rides with the forward
as a fixed function of moneyness.  (``To first order'' is doing real
work in this bullet --- \cref{sec:ssrboard} shows a case where the
leftover curvature term is worth volatility points.)
\item $\mathcal{R}=1$, \emph{sticky-strike}: the ATM volatility rises
to the value the old smile assigned to the new ATM strike --- each
strike kept its volatility, and the money walked down the skew to a
higher reading.
\item $\mathcal{R}=2$, \emph{sticky-local-vol}: the ATM volatility
rises \emph{twice} as far --- it overshoots the old curve.  The name,
and why two is the natural value for it, is the business of
\cref{sec:ssrhalf}.
\end{itemize}
A plain reading of the dial: $\mathcal{R}$ asks how much of the skew
is a forecast.  The skew says lower strikes carry higher volatility.
When spot actually goes lower, $\mathcal{R}$ is the fraction of that
extra volatility that the at-the-money reading actually
collects.

Compute the whole transport once by hand, on a toy smile where every
number is two-digit arithmetic.  Let today's smile be the straight
line $\sigma_{\mathrm{old}}(k)=0.20-0.15\,k$ (so $s_0=-0.15$, and the
ATM volatility is $20\%$), let $F=100$, and let the forward rise
$4\%$ ($H=+0.04$), as in \cref{sec:ssrsigns}.  (The taxonomy above
read the market \emph{downward}; the toy runs it \emph{upward}, so
every intuition from the bullets flips sign --- what rose there falls
here.)  Then
$\sigma_{\mathrm{old}}(H)=0.20-0.15\times0.04=0.194=19.4\%$ and
$s_0H=-0.15\times0.04=-0.006$, i.e.\ $-60$ vol bp, and
\eqref{eq:ssrshift} gives, at the new money
($k=0$) and at the old strike $K=100$ (new label $k=-H$, so the
re-indexed term is $\sigma_{\mathrm{old}}(0)$):
\[
\begin{array}{lcc}
 & \text{ATM (strike }104.08\text{)} & \text{strike }100\\[2pt]
\mathcal{R}=0\ \text{(sticky-moneyness)} & 20.0\% & 20.6\%\\
\mathcal{R}=1\ \text{(sticky-strike)}    & 19.4\% & 20.0\%\\
\mathcal{R}=2\ \text{(sticky-local-vol)} & 18.8\% & 19.4\%
\end{array}
\]
Check each against the two rules.  The ATM column is
$20\%+\mathcal{R}s_0H$: unchanged, down $60$ bp, down $120$ bp ---
\cref{prop:ssratm} exactly, since the toy smile has no curvature.  The
strike-$100$ column is $20\%+(\mathcal{R}-1)s_0H$: under
sticky-moneyness the strike is now below the money and picks up the
skew's $60$ bp; under sticky-strike it keeps $20\%$ by construction;
under sticky-local-vol it \emph{loses} $60$ bp, because the level term
pushes the whole curve past it.  The reader should reproduce all six
numbers on paper before going on; everything later in the chapter is
this computation with real curves.

\begin{figure}[htbp]
\centering
\paperfig{\textwidth}{fig_ssr_dial.pdf}
\caption{The dial (frozen hero smile).  (a)~The transport
\eqref{eq:ssrshift} built in its two moves for $\mathcal{R}=2$,
$H=-\MacSsrFanMovePct\%$: the dashed curve is the pure re-index
$\sigma_{\mathrm{old}}(k+H)$; the arrows add the uniform level
$(\mathcal{R}-1)s_0H=+\MacSsrDialLevelBp$ vol bp.  (b)~The exact ATM
response of the transport against the move for the three regimes
(solid), with the linear law $\mathcal{R}s_0H$ dotted: the slopes are
$\mathcal{R}s_0$, and all three curves share one bend --- at
$H=-6\%$ it is worth \MacSsrDialBendBp\ vol bp --- which is the
smile's own curvature entering through the re-index term.}
\label{fig:ssrdial}
\end{figure}

\Cref{fig:ssrdial} shows the same construction on the real node.  In
panel~(a), the orange curve is today's smile and the dashed grey curve
is the pure re-index $\sigma_{\mathrm{old}}(k+H)$ for a
$-\MacSsrFanMovePct\%$ move: the shape has slid bodily to the left,
and at the money it already reads higher --- that is the $s_0H$ from
the proof above.  The three arrows --- drawn at three different strikes to show the
lift is the same everywhere --- then add the level
$(\mathcal{R}-1)s_0H$, here $+\MacSsrDialLevelBp$ vol bp for
$\mathcal{R}=2$, lifting the dashed curve to the violet one.  Every
transported smile in this chapter is those two moves and nothing else.
Panel~(b) plots the exact ATM change $\sigma_{\mathrm{new}}(0)-
\sigma_{\mathrm{old}}(0)$ against the move for the three regimes, with
the straight lines $\mathcal{R}s_0H$ dotted on top.  Near zero move
the three solid curves are the three lines: slope $\mathcal{R}s_0$,
\cref{prop:ssratm}.  Away from zero they all bend by the same amount
--- \MacSsrDialBendBp\ vol bp at $H=-6\%$ --- and the bend has nothing
to do with $\mathcal{R}$: it is
$\tfrac12\sigma_{\mathrm{old}}''(0)H^{2}$, the smile's own curvature
picked up by the re-index term --- the \MacSsrDialBendBp\ vol bp is
that term measured on the fitted curve, and it is the first
$O(H^2)$ the proposition swept aside.  The same term explains a detail the sharp-eyed reader
may have caught in \cref{fig:ssrquestion}(a): the sticky-moneyness
tomorrow does not lie exactly on today's curve, though the words ``the
smile is a fixed function of moneyness'' say it should.  The transport
\eqref{eq:ssrshift} realizes each regime to first order in the move,
and its second-order error is exactly this one curvature term.

Two remarks close the section.  First, the level did not have to be
vertical: one could instead slide the argument farther,
$\sigma_{\mathrm{old}}(k+\mathcal{R}H)$, and match the same ATM
response.  The two versions differ by curvature times $H^{2}$ --- the
same second order the transport already declines to model --- so the
choice between them is a convention, and this book's is
\eqref{eq:ssrshift}.  Second, the dial is not unconstrained folklore:
within one-factor stochastic-volatility models --- models in which a
single random factor drives the volatility --- short-dated smiles
force $\mathcal{R}$ between the sticky-strike and sticky-local-vol
values \citep{Bergomi2016}, so the named regimes bracket the
modeling mainstream.  What no model choice can escape is
\cref{rem:ssrident}: today's surface prices every $\mathcal{R}$
identically.  The next section computes what the choice costs.
